% cover_letter.tex
% IEEE Transactions on Power Delivery — submission cover letter
% Paper: "Bayesian Three-Hypothesis SPRT for IBR-Aware Transformer
%         Differential Protection: Waveform Asymmetry, Fault-Integrity
%         Gating, and 30,000-Trial Monte Carlo Certification"
% Authors: A. V. Eluvathingal and K. S. Swarup, IIT Madras
\documentclass[11pt,a4paper]{article}
\usepackage[T1]{fontenc}
\usepackage[utf8]{inputenc}
\usepackage[margin=2.5cm]{geometry}
\usepackage{microtype}
\usepackage{parskip}
\usepackage{newtxtext}
\pagestyle{empty}

\begin{document}

\noindent
\today

\bigskip

\noindent
Editor-in-Chief\\
\textit{IEEE Transactions on Power Delivery}

\bigskip

\noindent
Dear Editor-in-Chief,

\medskip

We submit for your consideration the original research article:

\begin{center}
\textbf{``Bayesian Three-Hypothesis SPRT for IBR-Aware Transformer\\
Differential Protection: Waveform Asymmetry, Fault-Integrity Gating,\\
and 30,000-Trial Monte Carlo Certification''}\\[4pt]
\textit{Anoop V. Eluvathingal and K. Shanti Swarup}\\
Department of Electrical Engineering,
Indian Institute of Technology Madras, Chennai~600~036, India
\end{center}

\medskip\noindent\textbf{Motivation and problem addressed.}
Transformer differential (87T) protection has relied on second-harmonic (2H)
restraint to distinguish magnetising inrush from internal faults since the
1950s.
Inverter-based resources (IBRs) systematically destroy the assumption
underlying this restraint: the IBR current controller produces a sinusoidal
fault current with $I_2/I_1 < 2\,\%$, indistinguishable from inrush under
the 2H criterion.
Concurrently, residual-flux energisation ($B_r \geq 0.7\,$T) reduces the
inrush second-harmonic ratio to 8--10\,\%, below the blocking threshold,
creating a risk of false trips.
Monte Carlo evidence confirms that at $k_{\rm ibr} \geq 0.70$, the
conventional relay achieves a detection probability of only $P_D = 0.236$
--- one internal fault in four goes undetected.
No published method simultaneously addresses all three discriminands
(IBR internal fault, magnetising inrush, and CT-saturation-masked external
fault) within a single statistical framework with formal error-probability
bounds.

\medskip\noindent\textbf{Contributions.}
This paper proposes and validates a three-contribution IBR-aware 87T relay
architecture:

\begin{enumerate}
  \item \textbf{Three-hypothesis SPRT on waveform asymmetry (C6):}
        The per-half-cycle asymmetry index $A_k$ is derived from first
        principles and shown to be statistically separable across the
        three hypotheses --- inrush ($H_0$, $\mu_{A_k} = 0.62$), IBR
        fault ($H_1$, $\mu_{A_k} = 0.00$), and CT saturation
        ($H_2$, $\mu_{A_k} = -0.65$) --- with KL divergences of
        1.15 and 1.42 nats respectively.
        Two parallel log-likelihood-ratio accumulators with Wald boundaries
        achieve $P_D = 1.000$ and median trip time $t_{50} = 20\,$ms across
        30\,000 Monte Carlo trials.

  \item \textbf{Five-priority integration architecture (C7):}
        The SPRT is embedded into existing 87T relay logic through five
        sequentially evaluated priorities, with the SPRT trip signal gated
        by the fault-integrity index $f_{\rm int} \geq 0.55$ at Priority~4.
        This restores $P_D^{\rm ibr}$ from 0.236 to 1.000 across 20\,000
        MC trials and reduces the IBR-fault median trip time from
        $\geq 60\,$ms to 20\,ms, without hardware modification.

  \item \textbf{Fault-integrity index with formal separation proof (C8):}
        The scalar index $f_{\rm int} = 1 - \|\mathbf{r}_n\|/I_{\rm op}$,
        derived from the LM estimator residual, provides a formally proved
        0.34\,pu separation between the maximum inrush value
        ($f_{\rm int}^{H_0} \leq 0.21$) and the minimum internal-fault
        value ($f_{\rm int}^{H_1} \geq 0.55$), confirmed numerically across
        all 5\,000 Group~B (inrush) MC trials with zero overlap.
\end{enumerate}

\medskip\noindent\textbf{Validation.}
The chain is validated on 48 deterministic scenarios across three staged
studies (TR-57a analytical, TR-57b EMT-surrogate, TR-64 five-priority
integration) with zero misclassification, and on 30\,000 MC trials
(10\,000 Group~A IBR fault, 5\,000 Group~B inrush, 5\,000 Group~C
CT-saturation, 10\,000 Group~D SG fault).
All groups achieve rate $= 1.000$ with 95\,\% Wilson CI lower bound
$\geq 0.9992$, meeting IEC~60255-111 performance targets.

\medskip\noindent\textbf{Fit for \textit{IEEE Transactions on Power Delivery}.}
The paper addresses the transformer protection reliability gap created by
IBR integration~---~a problem directly relevant to the Transactions'
readership of protection engineers, utility planners, and relay manufacturers.
The combination of Wald's sequential hypothesis testing, formal distribution
derivation from physical models (Jiles--Atherton, Park, CT knee-point
saturation), and 30\,000-trial MC certification delivers the analytical and
empirical depth appropriate for a Transactions paper.
The work is original and has not been submitted elsewhere.

\medskip\noindent\textbf{Suggested reviewers.}
\begin{itemize}
  \item Prof.\ Ali Hooshyar, University of Toronto (IBR-impacted protection,
        transformer differential relaying)
  \item Prof.\ Bogdan Kasztenny, SEL Engineering (transformer inrush
        discrimination, differential relay design)
  \item Prof.\ Massimo Mitolo, Irvine Valley College / IEEE PSRC
        (power transformer protection, harmonic analysis)
  \item Dr.\ Evangelos Farantatos, EPRI (IBR fault characterisation and
        protection impacts)
\end{itemize}

\medskip\noindent\textbf{Authorship and conflicts.}
Both authors declare no conflicts of interest.
The contributions statement follows separately.
Simulation code and parametric data are available upon request.

\bigskip

\noindent Yours sincerely,

\bigskip\bigskip

\noindent
\textbf{Anoop V. Eluvathingal} \hfill \textbf{K. Shanti Swarup}\\
PhD Scholar (Corresponding Author) \hfill Professor\\
Department of Electrical Engineering \hfill Department of Electrical
Engineering\\
IIT Madras, Chennai 600 036, India \hfill IIT Madras, Chennai 600 036,
India\\
\texttt{ee21d017@smail.iitm.ac.in} \hfill
\texttt{swarup@ee.iitm.ac.in}

\end{document}
